\mode*
\section{The compile model}
\label{sec:compilation}

The first educational objective is the model that separates LaTeX from the
word processor: the document is \emph{source text}; the typeset result is
a distinct artifact produced by \emph{compilation}; what you see is not
what you edit.

\subsection{Documented difficulties}

The seed round found no studies of how learners (mis)understand the
compile model---the central gap this paper must fill (cf.\
\cref{sec:method-course-data}).
What the controlled comparison shows is consistent with a costly model
shift: users of the source-based system were slower and made more errors
on tasks where the WYSIWYG model suffices, while the source-based system
won on equations \autocite{Knauff2014Latex}.
The nearest learner evidence is from a neighbouring markup language:
students in a first web-development course \enquote{grappled with the
formal nature of [the] notation}, making \enquote{minute errors with case
sensitivity, white space, and missing delimiters}---a formality that
\enquote{contrasts with the flexibility not only of natural language, but
also of popular computing systems like word processors}
\autocite{Park2011WebDev}\footnote{Quoted from the dissertation reporting
the same study \autocite[p.~61]{Park2014openHTML}.}---and the same line of
work notes that WYSIWYG editors \enquote{can lead to a misapplication of
analogy, where learners intuit more than is warranted from the document
metaphor} \autocite[p.~38]{Park2014openHTML}.
% TODO(#2): Course-data collection (datintro LaTeX lab) to document the
% actual difficulties: e.g. editing the PDF-viewer window, recompiling
% never, or treating error output as terminal failure. These must be
% observed and coded, not assumed.

\subsection{Candidate critical aspects}
\label{sec:compilation-aspects}

Reading the difficulties above for the aspect the learner had not
discerned gives four candidates.
Because no study of learners and the compile model exists, we say for
each what it rests on: a documented difficulty in the literature, an
observation from our own course, or the structure of the object of
learning itself.

\begin{description}
  \item[Two artifacts, one direction]
    Edits happen in the source; the output is derived from it, never
    edited.
    No study documents the opposite conception for LaTeX; the aspect is
    posited from the object of learning and from the prior model the
    learners bring, whose document metaphor invites them to
    \enquote{intuit more than is warranted}
    \autocite[p.~38]{Park2014openHTML}, and whose cost when shifted the
    controlled comparison shows for professionals
    \autocite{Knauff2014Latex}.
    Tested by the item \emph{where the change goes} and, at the start of
    the module, by its open twin \emph{a mistake in the PDF}
    (\cref{app:knowledge-items,app:open-items}).
  \item[The compiler reads, it does not watch]
    Nothing changes until the document is compiled again; the PDF is a
    snapshot of the last run.
    Posited from the object of learning: it is the single-window word
    processor's live update that the learner has to unlearn, and no
    study documents how.
    Tested by the item \emph{after you edit} and its open twin \emph{the
    PDF has not changed}.
  \item[Whitespace and characters mean things]
    The source is a formal language, not decorated text: runs of spaces
    collapse, a blank line starts a paragraph, and a misplaced character
    is an error.
    Read off the finding that novices learning a markup language
    \enquote{grappled with the formal nature of [the] notation}, making
    \enquote{minute errors with case sensitivity, white space, and
    missing delimiters} \autocite{Park2011WebDev}; the difficulty is
    documented for HTML and CSS, and LaTeX's notation is at least as
    formal.
    Tested by the item \emph{spaces in the source}.
  \item[Structure nests]
    An environment opened with \verb|\begin| is closed by its \verb|\end|,
    a brace by its partner, and what lies between belongs to it.
    Read off the same study's finding that \enquote{hierarchies and
    paths, nesting} are among the concepts the novices' difficulties turn
    on \autocite{Park2011WebDev}, and off the existence of an instrument
    for reading hierarchies in code that counts LaTeX among its
    languages \autocite{Park2016Hierarchies}: the nested structure of the
    source is the one part of the compile model whose difficulty has been
    documented and measured.
    Tested by the item \emph{a list inside a list}.
\end{description}

Two of the four rest on difficulties documented for a neighbouring markup
language, none on a difficulty reported for LaTeX itself, because none
has been reported (\cref{app:protocol}); the open accounts of
\cref{sec:method-instruments} are where one would first show.

\subsection{Patterns of variation}

\begin{itemize}
  \item Contrast (system varies, document invariant): produce the same
    short document in a word processor and in LaTeX; locate where each
    stores \enquote{bold}.
  \item Contrast (source varies minimally): change one token, recompile,
    diff the output.
  \item Generalisation (model invariant, editor varies): the same
    source compiled from Overleaf and from the command line---the model,
    not the editor, is the invariant.
  \item Contrast (nesting varies, items invariant): the same three
    items as one flat list and as a list inside a list; the indentation
    and the labels show what belongs to what.
\end{itemize}
